\documentclass[11pt]{article}

\usepackage[margin=1in]{geometry}
\usepackage{parskip}
\usepackage{amsmath, amssymb, amsthm}
\usepackage{booktabs}
\usepackage{hyperref}

\title{A small library for solving real quadratic equations symbolically}
\author{Group N}
\date{}

\begin{document}
\maketitle

\begin{abstract}
We present \texttt{quadratic}, a small Python library for solving the
real quadratic equation \(a x^2 + b x + c = 0\). The library exposes
four short functions: the discriminant, the set of real solutions, the
number of real roots, and the coordinates of the vertex of the
corresponding parabola. The roots are returned exactly using
\texttt{sympy.solveset}~\cite{meurer2017}; the documentation follows
the Diataxis framework~\cite{procida2017} and the conventions of the
open-source textbook \emph{Python for Mathematics}~\cite{knight2024}.
\end{abstract}

\section{Introduction}

The real quadratic equation \(a x^2 + b x + c = 0\) is the first
non-trivial equation that mathematics students meet at school and one
of the few for which a closed-form solution is taught directly. The
solutions are given by the quadratic formula

\[
    x = \frac{-b \pm \sqrt{b^2 - 4 a c}}{2 a},
\]

and the nature of the roots depends on the sign of the discriminant
\(\Delta = b^2 - 4 a c\). When \(\Delta > 0\) there are two distinct
real roots; when \(\Delta = 0\) the two roots coincide; and when
\(\Delta < 0\) there are no real roots.

Several Python libraries can solve quadratics. The
\texttt{numpy.roots} function works with numeric inputs and returns
floating-point approximations, which is convenient but loses the exact
form of the answer. The standard \texttt{math} library does not handle
quadratics at all. We take a different approach. By calling
\texttt{sympy.solveset}~\cite{meurer2017} with the domain restricted
to \texttt{sympy.S.Reals} we get an exact \texttt{sympy.Set} of real
roots, which is also the recommended interface in the \texttt{sympy}
documentation.

We present a four-function library,
\texttt{quadratic}, that returns exact real solutions, the
discriminant, the number of real roots, and the vertex of the
parabola. The library is small but covers the three discriminant cases
in its tests.

\section{The library}

The library lives in the single module \texttt{quadratic.py} and
follows the conventions of the algebra chapter of \emph{Python for
Mathematics}~\cite{knight2024}: every input is wrapped in
\texttt{sympy.S(\dots)} before any arithmetic, so that an integer
input is automatically promoted to \texttt{sympy.Integer} and an
expression input passes through unchanged.

\texttt{discriminant(quadratic\_coefficient, linear\_coefficient,
constant\_coefficient)} returns \(b^2 - 4 a c\).

\texttt{solutions(quadratic\_coefficient, linear\_coefficient,
constant\_coefficient, variable)} returns the set of real roots. We
pass the constructed expression to \texttt{sympy.solveset} with
\texttt{domain=sympy.S.Reals}; this returns an empty set when
\(\Delta < 0\), a one-element set when \(\Delta = 0\) and a two-element
set otherwise.

\texttt{number\_of\_real\_roots(quadratic\_coefficient,
linear\_coefficient, constant\_coefficient)} returns \(0\), \(1\), or
\(2\) according to the sign of the discriminant.

\texttt{vertex(quadratic\_coefficient,
linear\_coefficient, constant\_coefficient)} returns the coordinates
\((-b/(2a), c - b^2/(4a))\) of the vertex of the parabola
\(y = a x^2 + b x + c\).

\section{Worked examples}

For \(x^2 - 3 x + 2 = 0\) the discriminant is
\(1 > 0\), so we expect two real roots. The library returns
\(\{1, 2\}\), which we verified by hand.

For \(x^2 - 2 x + 1 = 0\) the discriminant
is zero. The library returns the singleton \(\{1\}\), and the vertex
of \(y = (x - 1)^2\) is at \((1, 0)\), as expected.

For \(x^2 + 1 = 0\) the discriminant is
\(-4 < 0\). The library returns the empty set; the vertex of
\(y = x^2 + 1\) is at \((0, 1)\), above the \(x\)-axis, which is
consistent with the parabola not crossing it.

\section{Discussion}

The \texttt{sympy} documentation
recommends \texttt{solveset} over the older \texttt{solve} because it
returns a \texttt{Set}, which gives a uniform way to talk about the
empty case (no real roots) and the singleton case (repeated root).
Working with sets is also easier when the caller wants to combine the
result with set operations such as union or membership.

The library handles only quadratics. Cubic and
higher-degree polynomials would need a different formula (cubics) or
no closed form at all (degree five and above by Abel-Ruffini). A
natural extension would be to fall back to \texttt{sympy.solveset} for
higher-degree polynomials and expose the same interface.

\section{Conclusion}

We have presented \texttt{quadratic}, a small library for solving real
quadratic equations symbolically. The library returns exact roots, the
number of real roots, the discriminant, and the vertex of the parabola.

\bibliographystyle{plain}
\bibliography{references}

\end{document}
